% BEGIN REV09_PELL_CATALAN
\subsection{Dilatation de Pell, moment de Catalan et retour orienté}
\label{corpus9:iv:pell}

Le moment paramétrique \(\mathcal C_2\) précédent conserve la
variable de profondeur ; nous écrivons
\(G_{\mathrm{Cat}}=\mathcal C_2(1)\).
Sa spécialisation suivante reçoit
l’échelle de la vis spectrale de Pell
\[
 V^{-1}\mathscr S V=\lambda_P i\,\mathscr S,\qquad
 \lambda_P=2+\sqrt3,\qquad \rho=\lambda_P^{-1}=2-\sqrt3.
\]
La loi de covariance est l’hypothèse opératoire de cette
spécialisation ; elle s’entend dans le domaine commun invariant
de \(\mathscr S\) et du transport inversible \(V\).
L’unité \(\lambda_P\) appartient à \(\mathbb Q(\sqrt3)\) ;
l’unité \(3+2\sqrt2\) de la réalisation de compression
appartient à \(\mathbb Q(\sqrt2)\) et conserve son lecteur propre.
Dans son orbite transversale scalaire,
\[
 z_n=z_0(\lambda_P i)^n,\qquad
 r_n=r_0\lambda_P^n,\qquad
 \theta_n=\theta_0+\frac{\pi_{\mathrm{HMT}}n}{2}.
\]
Pour \(z_0\ne0\), éliminer \(n\) donne la suspension
\[
 r(\theta)=r_0\exp\!\left(
   \frac{2\log\lambda_P}{\pi_{\mathrm{HMT}}}
   (\theta-\theta_0)\right).
\]
La branche inverse a pour multiplicateur \(-i\rho\). Cette
réalisation d’échelle reçoit la vis comme antécédent
opératoriel et conserve son orientation.

\begin{proposition}[Lecture spectrale et évaluation à l’échelle contractante
]
\label{corpus9:iv:pell-evaluacion}
Dans \(\mathcal H=\ell^2(\mathbb N_{>0})\), soient

\[
 N\delta_n=n\delta_n,\qquad U_4\delta_n=i^n\delta_n,\qquad
 Q_4=\frac{U_4-U_4^*}{2i},\quad
 C_4=\frac{U_4+U_4^*}{2}.
\]
Pour \(0<r<1\) et \(\nu\in\mathbb C\), les traces
\[
 F_\nu(r)=\operatorname{Tr}(N^{-\nu}r^NQ_4),\qquad
 H_\nu(r)=\operatorname{Tr}(N^{-\nu}r^NC_4)
\]
convergent absolument et localement uniformément dans
\((\nu,r)\). On a \(F_2(r)=\mathcal C_2(r)\) et
\begin{align}
 \mathcal C_2(\rho)
   &=\frac23G_{\mathrm{Cat}}
     -\frac{\pi_{\mathrm{HMT}}}{12}\log\lambda_P,\\
 \mathcal D\mathcal C_2(\rho)&=\frac{\pi_{\mathrm{HMT}}}{12},
 &\mathcal D^2\mathcal C_2(\rho)&=\frac14 .
 \label{corpus9:iv:pell-valores}
\end{align}
\end{proposition}
\begin{proof}
La base diagonale donne les séries de coefficients
\(\sin(n\pi/2)\) et \(\cos(n\pi/2)\). Sur un compact de
\(\nu\) et avec \(r\leq r_0<1\), leurs valeurs absolues sont
dominées par \(n^M r_0^n\) pour un certain \(M\), série sommable.
La sélection des indices impairs dans la première trace
produit exactement \(\mathcal C_2(r)\).


La formule de soustraction de la tangente donne
\(\rho=\tan(\pi_{\mathrm{HMT}}/12)\), comme reconnaissance
angulaire ultérieure de l’unité de Pell. Pour
\(0<\theta<\pi/2\), l’intégration par parties du moment donne
\[
 \mathcal C_2(\tan\theta)
  =\theta\log(\tan\theta)-\int_0^\theta\log(\tan x)\,dx.
\]
Le terme de bord en zéro est nul. En soustrayant les séries
de Fourier de \(\log(2\sin x)\) et \(\log(2\cos x)\), et en
intégrant par régularisation d’Abel, il vient
\[
 -\int_0^\theta\log(\tan x)\,dx
   =\sum_{\substack{n\geq1\\2\nmid n}}
          \frac{\sin(2n\theta)}{n^2}.
\]
La régularisation admet une limite dans l’intégrale car la
singularité logarithmique en zéro est intégrable ; la série
intégrée est dominée par \(\sum n^{-2}\).

Pour \(n\) impair, soit \(\chi_{-4}(n)=(-1)^{(n-1)/2}\).
L’identité, vérifiée par les six classes impaires modulo
douze, est
\[
 2\sin(n\pi_{\mathrm{HMT}}/6)
   =\chi_{-4}(n)
     +3\,\mathbf1_{3\mid n}\chi_{-4}(n/3).
\]
Leur somme de poids \(n^{-2}\) donne
\(\frac12(1+3/9)G_{\mathrm{Cat}}=2G_{\mathrm{Cat}}/3\).
En prenant \(\theta=\pi_{\mathrm{HMT}}/12\) et
\(\log\rho=-\log\lambda_P\), on obtient la première
évaluation. Les autres découlent de
\(\mathcal D\mathcal C_2(r)=\arctan r\),
\(\mathcal D^2\mathcal C_2(r)=r/(1+r^2)\) et
\(\rho+\rho^{-1}=4\).

\end{proof}

La composante réelle de la même trace conserve
\[
 H_2(r)=\frac14\sum_{k\geq1}\frac{(-r^2)^k}{k^2},
 \qquad
 \mathcal D H_2(r)=-\frac12\log(1+r^2),\qquad
 \mathcal D^2H_2(r)=-\frac{r^2}{1+r^2}.
\]
Ces formules s’obtiennent en sélectionnant les indices pairs et
en dérivant la série absolument convergente. Par conséquent,
\[
 \operatorname{Tr}(N^{-2}r^NU_4)=H_2(r)+i\mathcal C_2(r).
\]
La dénomination
\(\operatorname{Li}_2(z)=\sum_{n\geq1}z^n/n^2\)
reconnaît ultérieurement cette trace comme
\(\operatorname{Li}_2(ir)\). En particulier,
\[
 \operatorname{Tr}(N^{-2}\rho^NU_4)
 =\frac14\operatorname{Li}_2(-\rho^2)
  +i\left(\frac23G_{\mathrm{Cat}}
       -\frac{\pi_{\mathrm{HMT}}}{12}\log\lambda_P\right).
\]
Le remplacement de \(U_4\) par \(U_4^*\) conserve la partie
réelle et inverse la partie imaginaire.

\begin{proposition}[Composition de profondeur et d’orientation
]
\label{corpus9:iv:pell-composicion}
Pour \(0<r,s<1\), \(a,b\in\mathbb Z/4\mathbb Z\) et \(m\geq0\), soit
\(B(r,a)=r^NU_4^a\). Alors
\[
 B(r,a)B(s,b)=B(rs,a+b),\qquad
 B(r,1)^m=r^{mN}U_4^m.
\]
En particulier, \(B(r,1)^4=r^{4N}\ne I\) pour \(0<r<1\).

\end{proposition}
\begin{proof}
Les deux facteurs sont diagonaux dans la même base. Dans
\(\delta_n\), leur produit multiplie par
\(r^ni^{an}s^ni^{bn}=(rs)^ni^{(a+b)n}\).
La formule des puissances en découle par récurrence, et
\(U_4^4=I\) donne le quatrième retour. Le facteur
\(r^{4N}\) reste distinct de l’identité.
\end{proof}

Le lecteur des puissances du multiplicateur inverse
\(-i\rho\) est \(B(\rho,-1)\). Son moment imaginaire a une
orientation opposée à celui de \(B(\rho,1)\), tandis que la
profondeur amortie est la même. Le transport \(V\),
l’observable \(\mathscr S\) et le lecteur \(U_4\) conservent
leurs espaces respectifs : la composition précédente se rapporte
aux multiplicateurs de l’orbite et à leur lecture diagonale.
La phase revient en quatre pas et la contraction accumulée
reste enregistrée.

% END REV09_PELL_CATALAN
